\section{Expected values of common distributions}\label{sec:distributions}

This section follows the article's table.  Finite distributions are exact
\(\ipra\) theorems.  Infinite and continuous distributions are effective
schemata whose parameters are rational or regular-real names with the displayed
positivity certificates.

\begin{proposition}[Bernoulli]\label{prop:bernoulli}
If \(X\in\{0,1\}\) with \(\Pp(X=1)=p\), then
\[
 \E[X]=p.
\]
\end{proposition}

\begin{proof}
The two weighted terms are \(0(1-p)\) and \(1p\).  In a natural-weight
presentation with weights \(W-r,r\), the equality is
\((0(W-r)+r)/W=r/W\).
\end{proof}

\begin{proposition}[Binomial]\label{prop:binomial}
For \(X\sim\operatorname{Bin}(n,p)\),
\[
 \E[X]=np.
\]
\end{proposition}

\begin{proof}
Construct \(X\) as the sum of \(n\) independent Bernoulli coordinates and use
finite linearity; independence is not actually needed for the mean once the
marginals are fixed.  Directly from the mass function,
\[
 \sum_{k=0}^{n}k\binom nkp^k(1-p)^{n-k}
 =np\sum_{k=1}^{n}\binom{n-1}{k-1}p^{k-1}(1-p)^{n-k}=np,
\]
using \(k\binom nk=n\binom{n-1}{k-1}\) and the finite binomial theorem.  For
rational \(p\), clear denominators; both combinatorial identities are proved by
primitive recursion.
\end{proof}

\begin{proposition}[Poisson]\label{prop:poisson}
For an effective nonnegative parameter \(\lambda\), the Poisson law
\[
 p_k=e^{-\lambda}\frac{\lambda^k}{k!}
\]
has expectation \(\lambda\).
\end{proposition}

\begin{proof}
At finite stage,
\[
 \sum_{k=0}^{N}kp_k
 =\lambda e^{-\lambda}\sum_{j=0}^{N-1}\frac{\lambda^j}{j!}.
\]
This is the exact shift identity \(k/k!=1/(k-1)!\).  Taylor's theorem with a
rational remainder bound gives a regular name for \(e^\lambda\) and proves
\(e^{-\lambda}\sum_{j\ge0}\lambda^j/j!=1\).  For a rational upper bound
\(L\ge\lambda\), choose \(N>2L\).  Successive absolute tail terms then have
ratio at most \(1/2\), so the tail is bounded by twice its first term.  This is
a primitive-recursive modulus.  Passing to the regular limit yields
\(\E[X]=\lambda\).
\end{proof}

\begin{proposition}[Geometric]\label{prop:geometric}
For rational \(0<p\le1\), let
\(\Pp(X=k)=p(1-p)^{k-1}\), \(k\ge1\).  Then
\[
 \E[X]=\frac1p.
\]
\end{proposition}

\begin{proof}
Put \(q=1-p\).  The finite polynomial identity
\[
 \sum_{k=1}^{N}kq^{k-1}
 =\frac{1-(N+1)q^N+Nq^{N+1}}{(1-q)^2}
\]
is proved by multiplying the sum by \(1-q\) twice and telescoping, with no
differentiation needed.  Multiplication by \(p=1-q\) gives a partial
expectation tending to \(1/p\).  Since \(q<1\), a rational separation
\(q\le1-1/d\) is computable from the rational code for \(p>0\).  The
finite binomial inequality \((1-1/d)^d\le1/2\) gives
\(q^{dm}\le2^{-m}\).  Bounded search therefore chooses \(m\) with
\((dm+d)2^{-m}\le2^{-r}\), yielding a primitive-recursive stage
\(N=dm\) at which \((N+1)q^N\le2^{-r}\).  This is an explicit tail
modulus; the case \(p=1\) has \(q=0\) and is immediate.
\end{proof}

\begin{proposition}[Uniform]\label{prop:uniform}
For rational \(a<b\), the uniform law on \([a,b]\) has
\[
 \E[X]=\frac{a+b}{2}.
\]
\end{proposition}

\begin{proof}
By exact polynomial integration,
\[
 \frac1{b-a}\int_a^b x\,dx
 =\frac{b^2-a^2}{2(b-a)}=\frac{a+b}{2}.
\]
Alternatively, the transformation \(x\mapsto a+b-x\) preserves the uniform
quadrature and pairs values with sum \(a+b\), forcing the mean to be its half.
The polynomial integral is already exact on every rational mesh.
\end{proof}

\begin{proposition}[Exponential]\label{prop:exponential}
For an effective positive parameter \(\lambda\), the density
\(f(x)=\lambda e^{-\lambda x}\) on \(x\ge0\) has
\[
 \E[X]=\frac1\lambda.
\]
\end{proposition}

\begin{proof}
On \([0,R]\), finite integration by parts gives
\[
 \int_0^R\lambda xe^{-\lambda x}\,dx
 =\frac1\lambda-e^{-\lambda R}\left(R+\frac1\lambda\right).
\]
The exponential Taylor bounds provide compact quadrature.  A positive rational
lower bound \(c\le\lambda\) gives
\(e^{-\lambda R}\le e^{-cR}\), and a primitive-recursive choice of \(R(p)\)
forces the displayed remainder below \(2^{-p}\).  Thus the improper integral
has regular value \(1/\lambda\).
\end{proof}

\begin{proposition}[Normal and standard normal]\label{prop:normal}
For an effective normal density with parameters \(\mu\) and \(\sigma>0\),
\[
 \E[X]=\mu.
\]
In particular, the standard normal mean is zero.
\end{proposition}

\begin{proof}
Substitute \(x=\mu+\sigma z\).  The expectation becomes
\[
 \mu\int_{-\infty}^{\infty}\varphi(z)\,dz
 +\sigma\int_{-\infty}^{\infty}z\varphi(z)\,dz,
\]
where \(\varphi(z)=(2\pi)^{-1/2}e^{-z^2/2}\).  The first integral is one by
the density certificate.  On every symmetric rational interval the second
integral is exactly zero because its integrand is odd and paired quadrature
nodes cancel.  A Gaussian tail certificate bounds the omitted absolute first
moment, so the improper limit is zero.  The transformation and all constants
are handled as regular interval names.
\end{proof}

\begin{proposition}[Pareto]\label{prop:pareto}
Let \(k>0\), \(\alpha>0\), and
\(f(x)=\alpha k^\alpha x^{-\alpha-1}\) for \(x\ge k\).  Then
\[
 \E[X]=
 \begin{cases}
 \dfrac{\alpha k}{\alpha-1},&\alpha>1,\\[4pt]
 +\infty,&0<\alpha\le1.
 \end{cases}
\]
Here rational powers are interpreted by regular names with monotonicity and
multiplication certificates.
\end{proposition}

\begin{proof}
The truncated mean is
\[
 \alpha k^\alpha\int_k^R x^{-\alpha}\,dx.
\]
For \(\alpha\ne1\), exact differentiation of the regular power name gives
\[
 \frac{\alpha k^\alpha}{1-\alpha}
   (R^{1-\alpha}-k^{1-\alpha}).
\]
If \(\alpha>1\), the first term tends effectively to zero and the remainder
simplifies to \(\alpha k/(\alpha-1)\).  If \(0<\alpha<1\), the positive term
\(R^{1-\alpha}\) has an explicit divergence modulus.  If \(\alpha=1\), the
truncated mean is \(k\log(R/k)\), which diverges by doubling \(R\).  All case
conditions are decidable for rational parameters; for regular parameters, a
strict separation certificate from one is required for the finite branch.
\end{proof}

\begin{proposition}[Cauchy]\label{prop:cauchy}
For a Cauchy law with location \(x_0\) and scale \(\gamma>0\), the expected
value is undefined.
\end{proposition}

\begin{proof}
Translation does not repair absolute integrability.  On the positive and
negative tails, \(\abs x\gamma/((x-x_0)^2+\gamma^2)\) is bounded below by a
positive multiple of \(1/\abs x\) beyond an explicitly computable radius.
Therefore both part expectations diverge logarithmically.  The signed
extended expectation is the indeterminate \(+\infty-\infty\), as in
\Cref{prop:cauchy-trunc}.
\end{proof}

\subsection{A mixed atom--density model}

\begin{proposition}[Precipitation occurrence and amount]
\label{prop:precipitation}
Suppose rain occurs with probability \(\alpha\), and conditional on rain the
amount is exponential with rate \(\lambda>0\); otherwise the amount is zero.
Then
\[
 \E[X]=\frac{\alpha}{\lambda}.
\]
\end{proposition}

\begin{proof}
Represent the model as a two-cell finite partition: dry with value zero and
wet with an effective exponential law.  The law of total expectation gives
\[
 \E[X]=(1-\alpha)\cdot0
 +\alpha\E[X\mid\text{wet}]
 =\alpha/\lambda.
\]
This is a direct combination of exact finite conditioning and the effective
exponential computation.
\end{proof}
